\documentclass[a4paper]{article}

\usepackage{amsmath}
\usepackage{amssymb}
\usepackage{graphicx} 
\usepackage{float}
\usepackage{caption}


\begin{document}
\section{Abstract}


\section{Introduction}

\section{Methods}

\subsection{Ensuring Network Connectivity}

Graph connectivity is necessary for this study, as it ensures that a
message originating at any node $s$ can reach any destination $t$.
We propose a simple heuristic to ensure connectivity.
The procedure repeatedly merges connected components, preserving the
vertex set and preferring degree-preserving edge swaps to limit the
addition of new edges.

For each edge swap, the procedure randomly selects a component $C_a$
containing a cycle and another component $C_b$ containing at least one
edge, giving preference to trees. A non-bridge edge $(u,v)$ from $C_a$
and any edge $(x,y)$ from $C_b$ are selected at random and replaced
with $(u,x)$ and $(v,y)$.

Since $C_a$ remains connected after its non-bridge edge is removed,
both endpoints of the removed edge in $C_b$ reconnect through $C_a$, even if that
edge was a bridge. The swap therefore merges the two components
while preserving every vertex degree and the total number of edges.

When no swap is possible, the procedure randomly selects two distinct
components and one vertex from each, then inserts an edge between them. This
increases the degree of each endpoint by one. Each swap or insertion
reduces the number of components by one, and the procedure continues
until the graph is connected.

\subsection{Edge capacity update heuristic}
We assume that the network operates in a non-critical regime,
where the message generation rate is sufficiently low to
allow the system to reach a statistical steady state.
At equilibrium, each edge $e \in E$ is associated with
a random variable $X_e$ representing its queue length.
We denote the cumulative distribution function of $X_e$ by $F_e$.

The free-flow rate (FFR) of an edge is defined as the
proportion of time during which its queue length does
not exceed its capacity. Formally,

\begin{equation}
    \mathrm{FFR}_e
    = \mathbb{P}(X_e \leq C_e)
    = F_e(C_e),
\end{equation}

\noindent
where $C_e$ denotes the capacity of edge $e$.

Given a desired minimum free-flow rate $\eta$, our
objective is to find a capacity configuration that
minimizes the total network capacity while ensuring
that every edge satisfies the prescribed FFR constraint.
This leads to the following optimization problem:

\begin{equation}
    \label{eq:optimization_problem}
    \begin{aligned}
        \underset{\mathbf{C} \in \mathbb{N}_+^{E}}
        {\text{minimize}}
        \quad & \sum_{e \in E} C_e \\
        \text{subject to}
        \quad & F_e(C_e) \geq \eta,
        \quad \forall e \in E.
    \end{aligned}
\end{equation}

\noindent
Here, $\mathbf{C} = (C_e)_{e \in E}$ denotes the
network capacity configuration.

To address this problem, we propose an iterative
heuristic based on local traffic observations. For each edge, we use
the observed queue lengths to estimate $F_e$ and select the smallest
positive capacity for which the estimated free-flow rate reaches
the target $\eta$. To reduce fluctuations between iterations, we
retain memory of previous observations when estimating $F_e$ and
adjust the capacity gradually toward the selected value.

Starting from an initial capacity configuration, we record the
queue length of each edge at every simulation step. The capacity
update method is applied every $\Delta T$ steps, with $k$ denoting its
iteration number. We denote by $\widehat{p}_e^{(k)}(q)$ the fraction of
observations in the current interval for which edge $e$ has queue
length $q$. This defines the empirical queue-length distribution.

To retain information from earlier intervals, we combine the current
histogram with the previous estimate. We denote this distribution
with memory by $\overline{p}_e^{(k)}$, which we update as follows:
\begin{equation}
    \overline{p}_e^{(k)}(q)
    = (1-\alpha)\widehat{p}_e^{(k)}(q) + \alpha \overline{p}_e^{(k-1)}(q) ,
\end{equation}
where $\alpha \in [0,1)$ is the momentum parameter. At the first
iteration, we initialize $\overline{p}_e^{(1)} = \widehat{p}_e^{(1)}$.

We estimate capacity from the distribution with memory
$\overline{p}_e^{(k)}$, rather than directly from the current empirical
distribution $\widehat{p}_e^{(k)}$. For $\alpha>0$, this incorporates
observations from previous iterations and reduces the influence of
fluctuations within a single sampling interval on the capacity estimate.
The cumulative distribution used for this estimate is therefore
\begin{equation}
    \widehat{F}_e^{(k)}(q) = \sum_{j=0}^{q} \overline{p}_e^{(k)}(j).
\end{equation}
We denote by $\widehat{C}_e^{(k)}$ the estimated capacity needed to satisfy
the target $\eta$ at iteration $k$. It is the smallest positive integer
capacity at which the estimated cumulative probability reaches $\eta$:
\begin{equation}
    \label{eq:minimum_capacity}
    \widehat{C}_e^{(k)} = \max\left(1,
        \min\left\{q \in \mathbb{N}_0 \;\middle|\;
        \widehat{F}_e^{(k)}(q) \geq \eta\right\}\right).
\end{equation}

The capacity is then updated using the geometric mean of its current
value and $\widehat{C}_e^{(k)}$. Since capacities are integers, we round this
mean upward to an integer:
\begin{equation}
    \label{eq:capacity_update}
    C_e^{(k+1)}
    = \Bigl\lceil \bigl(C_e^{(k)} \widehat{C}_e^{(k)}\bigr)^{1/2} \Bigr\rceil.
\end{equation}
This makes the update more gradual than directly assigning the
estimated capacity. The observations from the current interval are
then cleared, while $\overline{p}_e^{(k)}$ is retained for the next update.

The method has no automatic stopping mechanism. The number of updates
is specified in advance. In our simulations, we assessed whether this
number was sufficient by checking
that the mean edge capacity had stabilized and that the free-flow rate
target $\eta$ was satisfied for every edge. This method is a heuristic:
we have not proved convergence or global optimality. Nevertheless, we
obtained good results with this heuristic in our simulations.

\section{Results}

\section{Discussion}

\section{Conclusions}

\bibliographystyle{plain}
\bibliography{references}
\end{document}
